\section{The exact cubic threshold}\label{sec:threshold}

We now study $a_k(N)$ and $R_k=a_k/L_k$ from \eqref{eq:leading}, always with $2\le k\le N$. The matrix size is denoted by $N$ to distinguish this analysis from the depleted permutation in Section~\ref{sec:mixed}.

\begin{theorem}[Necessary and sufficient scale]\label{thm:threshold}
For every $2\le k\le N$,
\begin{equation}\label{eq:lower}
R_k-1\ge\frac{2k^3}{27N}.
\end{equation}
For $2k\le N$,
\begin{equation}\label{eq:upper}
R_k\le\exp(3k^3/N).
\end{equation}
Consequently, for every integer sequence $2\le k_N\le N$ as $N\to\infty$,
\begin{align}
R_{k_N}\to1&\quad\Longleftrightarrow\quad k_N^3/N\to0,\label{eq:iff}\\
R_{k_N}=O(1)&\quad\Longleftrightarrow\quad k_N^3/N=O(1).\label{eq:bounded}
\end{align}
The lower-bound constant $2/27$ is optimal uniformly over both parities.
\end{theorem}

The necessity in this result addresses the question in Collins and Nagatsu's published Remark 3.4 \cite{CN}. The statement concerns relative accuracy of their specific leading terms. It does not say that centered products cease to have useful upper bounds outside the cubic range.

\subsection{The known recurrence and its normalization}

Lemma~\ref{lem:depleted} with $r=0$, after rescaling, gives
\begin{equation}\label{eq:momentrec}
a_k=\frac{2(k-1)}{N(N-k+1)}a_{k-1}
+\frac{k-1}{N^2(N-k+1)}a_{k-2},
\qquad a_0=1,\quad a_1=0.
\end{equation}
This is \cite[equation (16)]{CN}. Its coefficients show that $a_k>0$ for $2\le k\le N$.

Set $E_k=R_k-1$. Substitution of \eqref{eq:leading} into \eqref{eq:momentrec} gives the following normalized recurrences, also present in the proof of \cite[Proposition 3.3]{CN}. For $m\ge2$,
\begin{align}
E_{2m}={}&\frac{N}{N-2m+1}E_{2m-2}
+\frac{4(2m-1)(2m-2)}{3(N-2m+1)}E_{2m-1}\notag\\
&+\frac{(2m-1)(8m-5)}{3(N-2m+1)}.\label{eq:evenerror}
\end{align}
For $m\ge1$,
\begin{align}
E_{2m+1}={}&\frac{3N}{(N-2m)(2m+1)}E_{2m}
+\frac{(2m-2)N}{(N-2m)(2m+1)}E_{2m-1}
+\frac{2m}{N-2m},\label{eq:odderror}
\end{align}
where the term involving $E_1$ is omitted at $m=1$. The initial error is $E_2=1/(N-1)$. Every coefficient and forcing term is nonnegative. In particular, $E_k\ge0$.

\subsection{Accumulating a lower bound}

In \eqref{eq:evenerror}, discard the second nonnegative term and bound the other denominator by $N$. This gives
\[
E_{2m}\ge E_{2m-2}+\frac{(2m-1)(8m-5)}{3N}.
\]
At $m=1$, the same inequality means $E_2\ge1/N$, with $E_0=0$ used only for the summation. Summing yields
\begin{equation}\label{eq:evenlower}
E_{2m}\ge\frac{C_{\rm e}(m)}N,
\qquad C_{\rm e}(m)=\frac{m(16m^2-3m-4)}9.
\end{equation}
The right side is at least $m^3/N$, because
\[
C_{\rm e}(m)-m^3=\frac{m(m-1)(7m+4)}9\ge0.
\]

For the odd orders define
\[
C_{\rm o}(m)=\frac{80m^3+195m^2+124m+6}{135}.
\]
A direct polynomial calculation gives $C_{\rm o}(1)=3$ and
\begin{equation}\label{eq:oddpoly}
C_{\rm o}(m)=\frac{3C_{\rm e}(m)+(2m-2)C_{\rm o}(m-1)}{2m+1}+2m.
\end{equation}
Equation~\eqref{eq:odderror} implies
\[
E_{2m+1}\ge\frac{3E_{2m}+(2m-2)E_{2m-1}}{2m+1}+\frac{2m}N.
\]
Induction using \eqref{eq:evenlower} and \eqref{eq:oddpoly} proves
\begin{equation}\label{eq:oddlower}
E_{2m+1}\ge C_{\rm o}(m)/N.
\end{equation}
At $m=1$, this starts from $E_3\ge E_2+2/N\ge3/N$. Finally,
\[
C_{\rm o}(m)-\frac{2(2m+1)^3}{27}
=\frac{75m^2+64m-4}{135}>0.
\]
For the even orders, $m^3\ge2(2m)^3/27$. Together these inequalities prove \eqref{eq:lower}.

\subsection{A self-contained upper bound}

Suppose $2k\le N$. In the recurrences for $R_j$, obtained by adding one back to \eqref{eq:evenerror} and \eqref{eq:odderror}, the sum of the coefficients is
\[
1+\frac{(2m-1)(8m-5)}{3(N-2m+1)}
\quad\text{for }j=2m,
\]
and $N/(N-2m)$ for $j=2m+1$. In both cases it is at most $1+3j^2/N$, using $N-j+1\ge N/2$. The same bound starts the induction at $R_2=N/(N-1)$. Positivity gives
\[
\max(1,R_2,\ldots,R_k)
\le\prod_{j=2}^k(1+3j^2/N)
\le\exp\left(\frac3N\sum_{j=2}^k j^2\right)
\le\exp(3k^3/N).
\]
This proves \eqref{eq:upper}.

If $k_N^3/N\to0$ or merely stays bounded, then $2k_N\le N$ eventually. The two bounds now imply \eqref{eq:iff} and \eqref{eq:bounded}. Sharpness of $2/27$ follows from Corollary~\ref{cor:small} below.

\begin{example}[Order four]
The recurrence gives
\[
a_4(N)=\frac{3(N+6)}{N^4(N-1)(N-2)(N-3)},\qquad
R_4(N)=\frac{N^2(N+6)}{(N-1)(N-2)(N-3)}.
\]
Thus $R_4=1+12/N+O(N^{-2})$. A fixed-order expansion like this is useful for checking normalization. The threshold theorem needs the uniform argument above, because its order is allowed to grow with $N$.
\end{example}
